\begin{solapartat}
\punts{0,25} $E = \frac{1}{2}kx^2$

\punts{0,4} Es pot calcular la constant elàstica a partir de qualsevol punt de la corba, en particular busquem un punt allunyat de l'origen per tenir més resolució:

Per $x = 15\ \mathrm{cm}$, $E = 2\ \mathrm{J} \Longrightarrow k = \dfrac{2E}{x^2} = \dfrac{2\cdot 2,0}{0,15^2} = 178\ \mathrm{N/m}$

\punts{0,1} $T = \dfrac{6,52}{10} = 0,652\ \mathrm{s}$

\punts{0,1} $\omega = \dfrac{2\pi}{T} = 9,64\ \mathrm{rad/s}$

\punts{0,4} $\omega = \sqrt{\dfrac{k}{m}} \Longrightarrow m = \dfrac{k}{\omega^2} = \dfrac{178}{9,64^2} = 1,91\ \mathrm{kg}$
\end{solapartat}

\begin{solapartat}
\figura[0.5]{sol_fig1.png}

\punts{0,25} Cal que s'identifiqui en el gràfic l'energia total del sistema que està al voltant de 0,9~J

\punts{0,5} Cal que s'estableixi que l'energia mecànica és constant.

\punts{0,5} Cal que s'identifiqui que la suma de les energies cinètica i potencial és igual a l'energia mecànica.
\end{solapartat}
